\section{Related work}\label{sec:related}

\paragraph{Automated physical model building} AutoPMoB collects documents about a target process,
extracts equations, variables, and experimental data, judges the equivalence of the extracted
information across documents, and reorganizes it into the desired model
\citep{AutoPMoBCSTR,AutoPMoBProto}. For extraction, an annotation tool and an algorithm collect
variable symbols from engineering documents \citep{VARAT}, and template sentences augment the data
for extracting variable definitions \citep{Nagayama}. For equivalence judgment, ProcessBERT, a
language model pre-trained on a chemical-engineering corpus, judges whether two variable
definitions refer to the same quantity \citep{ProcessBERT}, and a computer-algebra algorithm judges
whether two groups of equations are equivalent \citep{DAEEquiv,ZhangEquiv}. The work closest to
this study addresses the model-building step itself: \citet{KatoESCAPE} proposed an algorithm that
builds a model from extracted equations, and \citet{CombineEq} formulated model building as the
combination of extracted equations into models that satisfy four predefined requirements and
proposed an algorithm that finds every such model. This study differs in its input and output. The
combination method takes equations and their variables and returns every combination that meets the
requirements. This study also takes a description of the desired model, reads a context text for
each equation, and ranks the equations that the description asks for above the rest. Finding the
equations that concern the requested model does not by itself require context: in our benchmark,
whose queries take their variables from the symbols of the ground-truth equations, variable overlap
decides most of the ranking (Section~\ref{sec:mechanism}). Context is needed where variables cannot
decide, as when several combinations meet the requirements and only the description tells which one
the user wants, and this study does not show that it helps there.

\paragraph{Equation discovery from numerical data} Symbolic regression searches for an equation
that fits numerical data, determining its form as well as its coefficients \citep{AIFeynman}.
\citet{SchmidtLipson} distilled conservation laws from motion-tracking data without prior knowledge
of physics, sparse identification of nonlinear dynamics keeps the fewest terms of the governing
equations that represent the data \citep{SINDy}, and ALAMO learns algebraic surrogate models that
are as accurate and as simple as possible from few simulations or experiments \citep{ALAMO}. The
difficulty lies in the search: the candidate forms span a very large combinatorial space, and
methods that learn from data alone leave unused the knowledge that scientists rely on. LLM-SR
brings that knowledge in by letting an LLM propose equation skeletons, whose parameters an
optimizer fits to the data \citep{LLMSR}. This study differs in where its equations come from: it
builds a model from equations collected from the literature, so it needs no data measured on the
target process, and it reads context, which numbers do not carry. The two approaches can work
together, because the retrieved equations, with their coefficients left free, can serve as
candidate skeletons for symbolic regression.

\paragraph{Information retrieval and reranking} Classical information retrieval ranks the items of
a collection by their relevance to a query \citep{Manning}: term-weighting schemes such as TF-IDF
score lexical overlap \citep{Salton}, BM25 refines such scores \citep{BM25}, and latent semantic
analysis absorbs vocabulary mismatch in a low-dimensional space \citep{LSA}. Modern systems add a
second stage, in which a learned model, today typically a pretrained language model, reranks the
candidates of a fast first-stage retriever \citep{BERT,BERTrerank}, and dense retrievers perform
the first stage with learned embeddings \citep{DPR}. Our method inherits this retrieve-then-rerank
design and trains its reranker with a pairwise ranking loss \citep{LTR}. These methods score each
item independently of the other retrieved items. The best-known departure, maximal marginal
relevance, penalizes a candidate for resembling the items already selected and thereby diversifies
the result list \citep{MMR}. Our task needs the opposite relation: a good candidate should connect
to the equations already selected, by sharing their variables, and supply the query variables they
still miss. Where maximal marginal relevance treats overlap as redundancy, equation-set retrieval
treats shared variables as the links that make the equations a system.

\paragraph{Generation by large language models} LLMs generate domain knowledge directly from a
prompt \citep{GPT3}, so asking an LLM to write out the equations of the desired model is a natural
alternative to retrieving them. Generation, however, carries no provenance: LLMs produce fluent
output that no source supports \citep{Hallucination}, and a user cannot trace a generated equation
to a source to check it. The check matters for physical models, whose equations are solved as one
system, so that a single wrong equation can corrupt every variable connected to it. An LLM can
instead select equations from a database, which keeps the link to the sources, and
Section~\ref{sec:llmresults} compares the proposed method with LLMs used in this way.

\paragraph{Mathematical information retrieval} Mathematical information retrieval searches for
formulas with queries that contain text, formulas, or both \citep{MathSurvey}. Its methods compare
formulas as strings of symbols, as operator trees whose paths are matched \citep{Approach0}, or as
embeddings of those trees \citep{TangentCFT}. For this study such methods are components, not
rivals: our second stage also compares equations, but it reads only the sets of variables that
they contain. Shared tasks such as NTCIR-12 MathIR \citep{NTCIR12} and ARQMath \citep{ARQMath}
retrieve single formulas ranked by similarity to the query. Our method presents equations meant to
be used together, so the value of one equation depends on the equations presented with it, and we
found no formula-search benchmark that imposes this joint condition.
